\section{Payment Endpoints}

Payment endpoints expose the mock payment flow used by the shop.

Payments are created automatically during checkout. The API does not expose a direct endpoint for creating payments manually. After checkout, the returned payment can be marked as successful or failed using the mock payment endpoints.

\subsection{Payment Response Shape}

Payment endpoints return the following response shape:

\begin{minted}[frame=single]{json}
{
    "id": "665000000000000000000040",
    "orderId": "665000000000000000000030",
    "userId": "665000000000000000000001",
    "provider": "mock",
    "status": "pending",
    "amount": 49998,
    "currency": "PLN",
    "createdAt": "2026-05-30T10:00:00.000Z",
    "updatedAt": "2026-05-30T10:00:00.000Z"
}
\end{minted}

Successful mock payments may also include:

\begin{minted}[frame=single]{json}
{
    "mockTransactionId": "550e8400-e29b-41d4-a716-446655440000",
    "paidAt": "2026-05-30T10:05:00.000Z"
}
\end{minted}

Failed mock payments may also include:

\begin{minted}[frame=single]{json}
{
    "failureReason": "Card rejected by mock provider",
    "failedAt": "2026-05-30T10:05:00.000Z"
}
\end{minted}

\subsection{Payment Provider Values}

The current implementation supports one payment provider:

\begin{itemize}
    \item \texttt{mock}.
\end{itemize}

\subsection{Payment Status Values}

Payments may use the following status values:

\begin{itemize}
    \item \texttt{pending},
    \item \texttt{paid},
    \item \texttt{failed}.
\end{itemize}

\subsection{Payment Access Rules}

Payment access is ownership-based.

\begin{itemize}
    \item A regular user can access only their own payment.
    \item An administrator can access any payment.
    \item A payment can be marked as successful or failed only while it has status \texttt{pending}.
\end{itemize}

\subsection{\texttt{POST /payments/:id/mock-success}}

\textbf{Access:} Authenticated user or administrator

\textbf{Description:} Marks a pending mock payment as paid.

\textbf{Path parameters:}

\begin{itemize}
    \item \texttt{id} --- payment ObjectId.
\end{itemize}

\textbf{Headers:}

\begin{minted}[frame=single]{http}
Authorization: Bearer <accessToken>
\end{minted}

\textbf{Request body:} None.

\textbf{Important rules:}

\begin{itemize}
    \item The payment must exist.
    \item Regular users can mark only their own payments as successful.
    \item Administrators can mark any payment as successful.
    \item The payment must have status \texttt{pending}.
    \item The payment status is changed to \texttt{paid}.
    \item A mock transaction identifier is generated.
    \item The related order is changed from \texttt{pending\_payment} to \texttt{paid}.
    \item The operation is transactional with the related order update.
\end{itemize}

\textbf{Example request:}

\begin{minted}[frame=single]{http}
POST /payments/665000000000000000000040/mock-success
Authorization: Bearer <accessToken>
\end{minted}

\pagebreak
\textbf{Example success response:}

\begin{minted}[frame=single]{json}
{
    "id": "665000000000000000000040",
    "orderId": "665000000000000000000030",
    "userId": "665000000000000000000001",
    "provider": "mock",
    "status": "paid",
    "amount": 49998,
    "currency": "PLN",
    "mockTransactionId": "550e8400-e29b-41d4-a716-446655440000",
    "createdAt": "2026-05-30T10:00:00.000Z",
    "updatedAt": "2026-05-30T10:05:00.000Z",
    "paidAt": "2026-05-30T10:05:00.000Z"
}
\end{minted}

\subsection{\texttt{POST /payments/:id/mock-failure}}

\textbf{Access:} Authenticated user or administrator

\textbf{Description:} Marks a pending mock payment as failed.

\textbf{Path parameters:}

\begin{itemize}
    \item \texttt{id} --- payment ObjectId.
\end{itemize}

\textbf{Headers:}

\begin{minted}[frame=single]{http}
Authorization: Bearer <accessToken>
Content-Type: application/json
\end{minted}

\textbf{Request body:}

\begin{minted}[frame=single]{json}
{
    "reason": "Card rejected by mock provider"
}
\end{minted}

\textbf{Body fields:}

\begin{itemize}
    \item \texttt{reason} --- optional string describing why the mock payment failed.
\end{itemize}

\textbf{Important rules:}

\begin{itemize}
    \item The payment must exist.
    \item Regular users can mark only their own payments as failed.
    \item Administrators can mark any payment as failed.
    \item The payment must have status \texttt{pending}.
    \item If \texttt{reason} is omitted or blank, the backend uses \texttt{Mock payment failed}.
    \item The payment status is changed to \texttt{failed}.
    \item The related order is changed from \texttt{pending\_payment} to \texttt{payment\_failed}.
    \item The operation is transactional with the related order update.
\end{itemize}

\textbf{Example request:}

\begin{minted}[frame=single]{http}
POST /payments/665000000000000000000040/mock-failure
Authorization: Bearer <accessToken>
Content-Type: application/json
\end{minted}

\begin{minted}[frame=single]{json}
{
    "reason": "Card rejected by mock provider"
}
\end{minted}

\textbf{Example success response:}

\begin{minted}[frame=single]{json}
{
    "id": "665000000000000000000040",
    "orderId": "665000000000000000000030",
    "userId": "665000000000000000000001",
    "provider": "mock",
    "status": "failed",
    "amount": 49998,
    "currency": "PLN",
    "failureReason": "Card rejected by mock provider",
    "createdAt": "2026-05-30T10:00:00.000Z",
    "updatedAt": "2026-05-30T10:05:00.000Z",
    "failedAt": "2026-05-30T10:05:00.000Z"
}
\end{minted}

\subsection{Common Payment Errors}

\textbf{Payment not found:}

\begin{minted}[frame=single]{json}
{
    "message": "Payment not found",
    "error": "Not Found",
    "statusCode": 404
}
\end{minted}

\textbf{Payment belongs to another user:}

\begin{minted}[frame=single]{json}
{
    "message": "You cannot access this payment",
    "error": "Forbidden",
    "statusCode": 403
}
\end{minted}

\pagebreak
\textbf{Payment is no longer pending:}

\begin{minted}[frame=single]{json}
{
    "message": "Payment is not pending",
    "error": "Conflict",
    "statusCode": 409
}
\end{minted}